\section[Conclusão]{Conclusão}

Durante o projeto Plan Line na AGES I, minha trajetória foi desafiadora e recompensadora. Desde o início, eu estava empolgado por participar do meu primeiro projeto AGES, mas também tinha consciência de que muitas das tecnologias envolvidas eram novas para mim. Enfrentar essa realidade não foi simples, porém a curva de aprendizado intensa acabou sendo uma das partes mais importantes da experiência.

Na Sprint 0, percebi que o projeto exigiria conhecimentos técnicos que eu ainda não dominava. Essa constatação foi desconfortável, mas também serviu como incentivo para estudar HTML, CSS, JavaScript e TypeScript. Minha disposição para aprender foi fundamental para superar os desafios iniciais e começar a contribuir com mais segurança.

Na Sprint 1, ao mudar meu foco para o frontend e trabalhar no menu lateral do projeto, consegui desenvolver habilidades práticas em React e Tailwind CSS. A conclusão antecipada dessa tarefa foi uma conquista pessoal, pois demonstrou minha capacidade de adaptação e reforçou minha confiança para lidar com tecnologias novas.

A Sprint 2 trouxe uma autocrítica importante. Minha contribuição foi menor do que eu gostaria, principalmente pela dificuldade com backend e pela sobrecarga de outras demandas acadêmicas. Esse período mostrou que entusiasmo e boa vontade não substituem organização de tempo e preparo técnico. Também evidenciou que, em projetos colaborativos, atrasos individuais e falta de planejamento impactam o andamento coletivo.

Na Sprint 3, dei um passo importante ao sair da zona de conforto e atuar no backend, trabalhando na criação de endpoints para CRUD. Mesmo com pouca experiência em Node.js, contei com apoio de colegas mais experientes e aprendi conceitos como métodos HTTP, Prisma, PostgreSQL e organização de APIs. Essa experiência foi decisiva para perceber minha afinidade maior com backend, algo que passou a influenciar minha visão sobre desenvolvimento de software.

Na sprint final, minha atuação foi mais voltada ao suporte aos colegas, especialmente na configuração de ambiente e banco de dados. Essa experiência reforçou que contribuir para um projeto não se resume a escrever código. Ensinar, orientar, destravar dificuldades e compartilhar conhecimento também são formas essenciais de participação em uma equipe.

Ao longo do projeto, percebi que a AGES se relaciona diretamente com as disciplinas do curso, pois transforma conceitos vistos em aula em situações reais de tomada de decisão, negociação, implementação e entrega. Temas como banco de dados, desenvolvimento web, engenharia de requisitos, arquitetura e métodos ágeis deixaram de ser apenas conteúdos teóricos e passaram a aparecer em problemas concretos do dia a dia do projeto.

Como ponto positivo, destaco o aprendizado técnico acelerado, o contato com stakeholders reais, a vivência em equipe e a possibilidade de experimentar diferentes frentes de desenvolvimento. Também considero muito valiosa a oportunidade de errar, revisar decisões e amadurecer a partir das dificuldades.

Como ponto a melhorar, reconheço que tanto eu quanto a equipe poderíamos ter administrado melhor o tempo, especialmente nas sprints intermediárias. Também poderíamos ter antecipado riscos, dividido responsabilidades de forma mais clara e evitado deixar tarefas críticas para perto do prazo final. Do ponto de vista pessoal, entendo que eu poderia ter sido mais proativo em alguns momentos, principalmente quando percebi dificuldades técnicas ou gargalos de organização.

Essa jornada não foi apenas sobre desenvolvimento de software. Foi também sobre desenvolvimento pessoal e profissional. Aprendi a lidar com desafios técnicos, a colaborar com colegas, a comunicar dificuldades, a ouvir stakeholders e a reconhecer a importância das relações humanas no ambiente de trabalho. A experiência na AGES I deixou uma marca importante na minha formação e foi decisiva para meu crescimento como desenvolvedor.
